\label{sec:related}

\subsection{基于互补约束的非光滑接触动力学}

非光滑接触动力学通常通过非穿透约束、速度层互补条件和库仑摩擦律建立模型。在多体系统中，这类方法具有较好的鲁棒性与工程可落地性，特别适合刚体碰撞、频繁接触和较大时间步长的仿真。Chrono 的 NSC 路线即属于该范式，其核心思想是将接触约束装配到系统描述器中，通过互补求解器或相关变体求解接触反力与系统速度更新。该类方法的主要优点是求解框架统一、适合复杂刚体系统耦合，但其接触层通常依赖离散接触点表示，因此对于连续区域接触的高精度建模存在天然张力。

\subsection{基于接触点缓存的流形维护}

在工程实现上，常见碰撞后端会在窄相检测之后维护少量持久接触点组成的接触流形，以提高时间连续性、暖启动效果和求解效率。该类方法的接触流形本质上更接近“接触点缓存”而非“连续接触区域的物理近似”。它们通常以启发式方式保留最深点、覆盖面积较大的点或时间上持续存在的点，因而在点接触和局部小面积接触中效果较好，但在长条接触、复杂多 patch 接触和大尺度连续区域接触中，可能出现接触作用线漂移、力矩失真与拓扑混叠。

\subsection{基于接触曲面和压力场的分布式接触}

与离散点流形不同，另一类方法从接触曲面出发，在连续接触区域上定义法向压力、切向牵引或相关势函数，再通过面积分得到结果力和结果力矩。此类方法更强调连续区域本身的几何与力学意义，因此在需要高精度 body-level wrench 或接触面演化分析时更具优势。然而，其代价往往是求解模型更复杂、与现有刚体互补求解框架的直接兼容性较弱。

\subsection{等效接触点与接触 wrench 表示}

等效接触点（ECP）相关工作指出，对单个凸接触区域，可以用一个等效接触点及其对应的接触 wrench 统一表征点接触、线接触与面接触。这一思想的重要启发在于：连续区域接触并不必然要求大量离散点才能在刚体动力学层面被精确表征；关键在于是否保留了正确的结果 wrench 与作用线信息。然而，直接使用 patch-level 六维接触元又需要改写现有求解器的数据结构与约束装配过程。

\subsection{SDF 几何与接触建模}

SDF 在接触建模中的优势主要体现在三个方面：其一，可直接提供局部 gap 与法向；其二，便于构造连续接触区域及其局部微分结构；其三，适合与体素、稀疏体数据结构和自适应采样耦合。现有基于 SDF 的接触研究多集中于穿深估计、法向计算、柔顺势函数构造和图形学仿真。相比之下，如何在 SDF 基础上面向刚体 NSC 框架构建“高保真但仍可低维求解”的接触约化，仍缺乏系统研究。

\subsection{本文方法与现有工作的区别}

本文方法既不同于传统少量缓存点的流形维护，也不同于直接在全局系统中引入 patch-level 六维接触元。本文的出发点是：以 SDF dense 接触区域为参考对象，在局部 patch / subpatch 层面保持其可行 wrench 锥和当前参考 wrench，并最终将这一等效结果重新投影回标准点接触集合。换言之，本文不是在求解器末端进行“删点”，而是在接触建模接口层进行“保持 wrench 等效的低维投影”。这使其既能显著改善连续区域接触的刚体层面等效精度，又能保持对 Chrono NSC 现有求解链路的兼容性。
